\documentclass{article}
\usepackage[utf8]{inputenc}
\usepackage{algorithm}
\usepackage{algpseudocode}
\usepackage{amsmath}

\begin{document}

% -----------------------------------------------------------------------------
% Resource-Aware Adaptive Partitioning Algorithm
% -----------------------------------------------------------------------------

\begin{algorithm}[hbt!]
    \caption{Resource-Aware Adaptive Granularity Partitioning}
    \label{alg:adaptive_partitioning}
    \begin{algorithmic}[1]
        % 输入输出定义
        \Require Top-level design hierarchy tree $T$, Resource threshold vector $R_{th}$ (LUTs, FFs, BRAMs)
        \Ensure Set of partitionable units $S$

        % 初始化
        \State $S \gets \emptyset$
        \State $Q \gets \text{Queue}()$
        \State $Q.\text{enqueue}(T.\text{root})$ \Comment{Start from the top-level module}

        % 主循环
        \While{$Q \text{ is not empty}$}
        \State $m \gets Q.\text{dequeue}()$
        \State $usage \gets \text{EstimateResource}(m)$ \Comment{Evaluate resource consumption}

        % 判断逻辑
        \If{$usage \le R_{th}$}
        \State $S \gets S \cup \{m\}$ \Comment{Module fits, preserve hierarchy}
        \Else
        \If{$\text{IsLeaf}(m)$}
        \State \Comment{Large leaf node (e.g., huge netlist block)}
        \State $sub\_blocks \gets \text{NetlistDecomposition}(m)$
        \State $S \gets S \cup sub\_blocks$
        \Else
        \State \Comment{Module too large, flatten one level}
        \State $children \gets \text{GetChildInstances}(m)$
        \For{$child \in children$}
        \State $Q.\text{enqueue}(child)$
        \EndFor
        \State \text{HandleGlueLogic}(m) \Comment{Add top-level glue logic to S}
        \EndIf
        \EndIf
        \EndWhile

        \Return $S$
    \end{algorithmic}
\end{algorithm}

\end{document}